\secciongrande{Lenguajes de Programación en Entorno Cliente: Dos Entornos de Ejecución}

\resumen{
  La interactividad en portales web requiere lenguajes específicos. Se distinguen entornos cliente (ejecutados en el navegador) y entornos servidor (ejecutados en el servidor). La elección del lenguaje depende de las necesidades del proyecto y de la habilidad del desarrollador.
}

\definicion{Entorno Cliente}{Conjunto de lenguajes ejecutados en el dispositivo del usuario, típicamente un navegador web, para interactuar con la interfaz.}
\definicion{Entorno Servidor}{Conjunto de lenguajes ejecutados en un servidor remoto para gestionar lógica de negocio y bases de datos.}

\subtituloA{Diferenciación y elección de lenguajes}

\lista{
  **Entorno Cliente:** Responsable de la interacción directa con el usuario, mejorando la atractividad de la web. |
  **Entorno Servidor:** Responsable de la lógica de negocio, gestión de datos y respuesta a peticiones del cliente.
}

\nota{Para el examen}{
Es crucial evaluar ventajas e inconvenientes de cada lenguaje; la decisión final recae en el desarrollador según su competencia.
}
